% Generated by `taskgen course examples` from ecc.point_counting (seed 3).
% Do not edit: change examples.yaml and regenerate.
\begin{example}[подсчёт точек с помощью символов Лежандра]\label{ex:ec-point-counting}
Проверить, задаёт ли уравнение \[y^2 = x^3 + 2x + 4\] эллиптическую кривую над $\mathbb{F}_{11}$. Если да, перечислить все точки $E(\mathbb{F}_{11})$ (с помощью символов Лежандра), найти $\#E(\mathbb{F}_{11})$ и проверить оценку Хассе; если нет, указать особую точку.

\par\noindent\textit{Решение.} $4a^3 + 27b^2 = 4 \cdot 2^3 + 27 \cdot 4^2 = 464 \equiv 2\pmod{11}$, т.\,е. $\Delta = -16(4a^3 + 27b^2) \ne 0$ и уравнение задаёт эллиптическую кривую.\par Квадраты в $\mathbb{F}_{11}$: $1^2 \equiv 1,\ \allowbreak 2^2 \equiv 4,\ \allowbreak 3^2 \equiv 9,\ \allowbreak 4^2 \equiv 5,\ \allowbreak 5^2 \equiv 3$; $\left(\frac{c}{11}\right) = 1$ для $c$ из этого списка и $-1$ для остальных $c \ne 0$. Для каждого $x$ вычисляем $f(x) = x^3 + 2x + 4$; при $\left(\frac{f(x)}{p}\right) = 1$ получаем две точки $(x, \pm y)$, при $f(x) = 0$ --- одну (точку порядка~$2$):\[\begin{array}{r|r|r|l}x & f(x) & \left(\frac{f(x)}{p}\right) & y \\ \hline 0 & 4 & 1 & 2,\ 9 \\ 1 & 7 & -1 & \text{---} \\ 2 & 5 & 1 & 4,\ 7 \\ 3 & 4 & 1 & 2,\ 9 \\ 4 & 10 & -1 & \text{---} \\ 5 & 7 & -1 & \text{---} \\ 6 & 1 & 1 & 1,\ 10 \\ 7 & 9 & 1 & 3,\ 8 \\ 8 & 4 & 1 & 2,\ 9 \\ 9 & 3 & 1 & 5,\ 6 \\ 10 & 1 & 1 & 1,\ 10 \\ \end{array}\]$E(\mathbb{F}_{11}) = \{\mathcal{O},\ (0, 2),\ \allowbreak (0, 9),\ \allowbreak (2, 4),\ \allowbreak (2, 7),\ \allowbreak (3, 2),\ \allowbreak (3, 9),\ \allowbreak (6, 1),\ \allowbreak (6, 10),\ \allowbreak (7, 3),\ \allowbreak (7, 8),\ \allowbreak (8, 2),\ \allowbreak (8, 9),\ \allowbreak (9, 5),\ \allowbreak (9, 6),\ \allowbreak (10, 1),\ \allowbreak (10, 10)\}$.\par $\#E = p + 1 + \sum_x \left(\frac{f(x)}{p}\right) = 11 + 1 + 5 = 17$. Оценка Хассе: $t = p + 1 - \#E = -5$, $\abs{t} = 5 \le 2\sqrt{11} \approx 6{,}63$ (равносильно $t^2 = 25 \le 4p = 44$) --- выполнена.\[\#E(\mathbb{F}_{11}) = 17.\]
\end{example}
